\section{A minimal physical and accounting model}\label{sec:modele}
\subsection{Building intuition with four entities}
Imagine four entities able to extract a resource from their
environment. Each has a mobilizable stock $K_i$: the larger it is,
the more it extracts, but each additional unit of stock yields less
than the previous one. The function $f_i(K)=A_iK^{\gamma_i}$
describes this \emph{extraction technology}. The input joins the
stock, a fraction of which decays at each step. New entities appear
with the same endowment $K^\circ$ and, in the reference regime,
the same production exponent $\gamma^\circ=1/2$.

A well-endowed entity can lend stock to another, where that stock
produces more at the margin. This real transfer simultaneously
creates a claim for the lender and a debt for the borrower. The
\emph{principal} $q$ is the nominal amount recorded; the \emph{interest
service} $rq$ is the payment due at each step. Interest moves the
resource between entities, without creating any. The principal is not
progressively repaid in this model: this is a strong institutional
assumption.

\begin{figure}[htbp]\centering
\begin{tikzpicture}[>=Stealth,font=\small,
 ent/.style={circle,draw=blue!60!black,fill=blue!7,minimum size=15mm,align=center},
 lab/.style={font=\scriptsize,fill=white,inner sep=2pt}]
\draw[rounded corners,fill=gray!3,draw=gray!50] (-1.5,-1.7) rectangle (9.5,2.3);
\node[ent] (a) at (0,0.7) {$e_1$\\$K_1=16$};
\node[ent] (b) at (4,0.7) {$e_2$\\$K_2=4$};
\node[ent] (c) at (7.5,1.2) {$e_3$\\$K_3=9$};
\node[ent] (d) at (7.5,-0.8) {$e_4$\\$K_4=1$};
\draw[->,very thick,blue!70!black] (a) to[bend left=18]
 node[lab,above] {initial loan $q=6$} (b);
\draw[->,dashed,thick,orange!80!black] (b) to[bend left=22]
 node[lab,below] {interest $rq$, at each step} (a);
\draw[->,green!45!black,thick] (-.9,2.8)--(a);
\node[lab] at (2.4,2.8) {input extracted from the environment: $f_i(K_i)$};
\draw[->,green!45!black,thick] (7.5,2.8)--(c);
\draw[->,red!70!black] (d)--(9.8,-1.4) node[right,font=\scriptsize] {erosion};
\node[align=center,font=\scriptsize] at (2.7,-1.25)
 {After the loan: $K_1=K_2=10$; $C_1=6$, $D_2=6$.\\
 The net worths remain $NW_1=16$ and $NW_2=4$.};
\end{tikzpicture}
\caption{\textbf{Real stock and accounting commitment are not the same thing.}
Fictitious example, inspired by the pedagogical diagram of the internship
presentation. The initial values are chosen to explain the mechanism, not
extracted from a simulation. Entities 3 and 4 show that not all entities
have to take part in the same exchange. Inputs and erosion are drawn on
only a few entities to lighten the diagram; the rules apply to all.
The external reservoir is not simulated.}
\label{fig:exemple}
\end{figure}

The loan thus preserves the instantaneous net worth of each partner,
but changes its productive capacity and its future commitments. An
entity may later be unable to honor its service, or owe more than it
holds. Its disappearance wipes out the claims held on it: its lenders
can in turn become insolvent. The coupling between real stocks,
persistent commitments and propagated losses is the mechanism studied
here.


\subsection{State, accounting, invariants}

The population is open and comprises a single type of agent, called an
\emph{entity}. No role -- lender, borrower, bank -- is declared;
all entities follow the same rules, and any observed differentiation
is emergent.

The variables $K_i,C_i,D_i,NW_i$, age and flows received or paid are
specific to entity $i$. The parameters $\lambda,\delta,\sigma,\rho$ and
the credit institution are common to the whole population.\footnote{The
parameters $\lambda$, $\delta$ and $\sigma$ are made explicit in
\S\ref{sec:pas}; $\rho$ and the credit institution in \S\ref{sec:credit}.
The glossary in appendix~\ref{sec:glossaire} gathers the notations.} The technology
$(A_i,\gamma_i)$ is attached to each entity; it is identical for all
at the reference point, but can differ after a targeted intervention.

An entity $i$ holds a real productive stock $K_i \ge 0$, counted in
joules, which we call its \emph{capital}. This is its intrinsic
productive size, not an accounting capital. A loan contract ties a
lender to a borrower through a nominal \emph{principal} $q > 0$ and a
per-step \emph{rate} $r > 0$. Summing over the contracts in which $i$
is involved gives its claims $C_i$ and its debts $D_i$, then its
\emph{net worth}
\begin{equation}
NW_i = K_i + C_i - D_i .
\label{eq:nw}
\end{equation}

The asymmetry between the two stocks is constitutive: capital is real,
productive and depreciable, whereas the principal is nominal, perpetual
and not amortized. The debt network can therefore remain heavy while
real assets erode.

The rules algebraically enforce, at every step, $K_i \ge 0$,
$\sum_i C_i = \sum_i D_i$, the equality between the sum of the net
worths of the survivors and their total capital, and the real balance
sheet
\begin{equation}
\sum_i K_i(t{+}1) - \sum_i K_i(t) = \text{endowments} + \text{shock}
  + \text{production} - \text{depreciation} - \text{destruction},
\label{eq:bilan}
\end{equation}
the five terms on the right being recorded separately at each step.
The numerical check accounts for floating-point rounding.
Loans and interest are purely internal transfers: there is neither
creation nor leakage of joules outside the channels named
in~\eqref{eq:bilan}.

\subsection{A time step}\label{sec:pas}

The phases run in this order, which is part of the model.
Scheduled interventions are applied before births.

\begin{enumerate}[label=(\roman*),leftmargin=2.2em]
\item \textbf{Births.} A number $B_t \sim \mathrm{Poisson}(\lambda)$
  of entities enters, each endowed with a capital $K^\circ$. The parameter $\lambda$ is
  the \emph{birth intensity} and $K^\circ$ the \emph{birth endowment}.
\item \textbf{Shock.} $K_i \leftarrow K_i \, e^{\eta_i}$, with
  $\eta_i \sim \mathcal{N}(-\sigma^2/2,\ \sigma^2)$ independent. The
  multiplicative mean equals one: no drift is imposed. The parameter
  $\sigma$ fixes the common amplitude of the noise; the draw $\eta_i$ is
  individual and renewed at each step. This is not a shared global
  shock.
\item \textbf{Production.} $K_i \leftarrow K_i + A\,K_i^{\gamma}$, with
  $0<\gamma<1$. The parameter $A$ is the \emph{productivity} of the
  technology and $\gamma$ its \emph{production exponent}. Concavity is
  essential: it is what creates a gain from pooling
  (\S\ref{sec:credit}).
\item \textbf{Interest service.} This is the periodic payment of the
  loan charge, without repayment of the principal. Each borrower pays
  $\sum r q$ drawn from its capital; if it cannot, it pays pro rata and
  enters service default.
\item \textbf{Depreciation.} $K_i \leftarrow (1-\delta)K_i$, where
  $\delta$ is the \emph{per-step depreciation rate}. It represents the
  temporal erosion of the stock; its value therefore depends on the
  duration chosen for a step (\S\ref{sec:temps}). Principals, for their
  part, do not depreciate.
\item \textbf{Market.} See \S\ref{sec:credit}.
\item \textbf{Bankruptcies.} See \S\ref{sec:faillite}.
\item \textbf{Measurements.} No logging option consumes a random draw:
  trajectories with and without measurement are identical.
\end{enumerate}

\subsection{Stocks, flows and the physics--economics correspondence}
\label{sec:physique}
The parallel between physics and economics is the center of the
approach, not an analogy added after the fact. It is nonetheless
necessary to distinguish a correspondence of dimensions, a
representational hypothesis and an empirical validation.
The model is a fundamental study of interactions: the entity remains
generic, and can shed light on a concrete system if the relevant
mechanisms are indeed present there, without implicitly becoming a
person and then a firm.

\begin{center}
\begin{tabularx}{\linewidth}{@{}lXX@{}}
\toprule
Object & Dimension and operation & Proposed economic reading \\
\midrule
$K_i$ & Stock, in model joules & Mobilizable productive size \\
$A_iK_i^{\gamma_i}$ & Production added during a step, in joules & Productive flow, not a wage \\
$A_iK_i^{\gamma_i}/\Delta t$ & Mean power over a step & Production capacity per unit of time \\
$q$ & Nominal principal, same unit of account as $K$ & Claim and debt, not amortized \\
$rq$ & Transfer during a step & Interest service, not resource creation \\
$NW_i$ & Accounting stock $K_i+C_i-D_i$ & Balance-sheet position, distinct from the physical stock \\
\bottomrule
\end{tabularx}
\end{center}
A watt is one joule per second. The engine's outputs are in joules
\emph{per step}; calling them watts would require fixing $\Delta t$ in
seconds, which has not been calibrated.\footnote{In the update
$K\leftarrow K+A K^\gamma$, $A$ has dimension
$\mathrm{J}^{1-\gamma}$ for the production of one step. The
coefficient of a corresponding differential equation would have
dimension $\mathrm{J}^{1-\gamma}/\mathrm{s}$. Comparing them without
specifying the step would make a physical dimension disappear. The
rate $r$ is a fraction due per step; $r/\Delta t$ has the inverse
dimension of a time.}

The interpretive hypothesis is that certain monetary transfers
represent transfers of entitlements to mobilize exergy. The real stock
and the nominal network are thus coupled, but not identical: a loan
moves $K$ and simultaneously records a claim and a debt. Production
represents an input to the modeled stock; the external reservoir that
enables it is not described. The internal balance sheet is closed in
accounting terms, not thermodynamically over the whole
economy--environment system. Furthermore, the observable called
"total income" is $\mathrm{prod}+\mathrm{int\_in}$: it adds
production and transfers received, and must not be equated without
correction to an aggregate net production.

\subsection{Scales: which parameters are truly independent?}
\label{sec:parametres_reduits}
In the homogeneous regime, without noise or credit, the order
production then depreciation gives
\[
 K_{t+1}=(K_t+A K_t^\gamma)(1-\delta),\qquad
 K_{\rm aut}=\left[\frac{A(1-\delta)}{\delta}\right]^{1/(1-\gamma)}
\]
for $A>0$, $0<\delta<1$ and $0<\gamma<1$. $K_{\rm aut}$ is the stock toward which an isolated entity, without
credit or noise, tends when extraction and depreciation offset each
other. This deterministic fixed point is not a mean of the noisy
system. Setting $x_i=K_i/K_{\rm aut}$ transforms
production into
\[
 x_i\longleftarrow x_i+\frac{\delta}{1-\delta}x_i^\gamma .
\]
At fixed institution, the common scale of $A$ and $K^\circ$ can thus
be replaced by the ratio $\kappa^\circ=K^\circ/K_{\rm aut}$: it
measures the size of a new entity relative to what it could reach
alone. A low value means it enters far below this scale. For a
toy model, avoiding double-counting the same scale freedom is
essential.\footnote{An equivalent option is to choose $K^\circ$ as the
unit and to retain $a=A (K^\circ)^{\gamma-1}$. The point here is to
eliminate \emph{one} unit-of-capital freedom, not to show that the
other parameters are redundant: $\gamma$, $\delta$, $\sigma$,
$\lambda$ and the institution remain relevant. An invariance of
certain statistics with respect to $\lambda$ is not a trajectorial
equivalence of populations of different sizes.}

The rescaling $K,q,K^\circ\mapsto c(K,q,K^\circ)$ and
$A\mapsto c^{1-\gamma}A$ preserves the homogeneous equations and the
marginal rates. Absolute numerical thresholds and the engine's
rounding limit software exactness; recalibrating the time step is not
an exact symmetry of the discrete market: the numerical study is
presented in \S\ref{sec:temps} and the non-equivalence of the
operations is detailed in appendix~\ref{sec:composition_temps}. A rise
in $A$ at fixed $K^\circ$ therefore changes $\kappa^\circ$, whereas its
compensation $K^\circ\mapsto K^\circ\,m^{1/(1-\gamma)}$ preserves it.
The two are distinct experiments: compensation is a control, not an
imposed physical correction.

The tension $T=K_{\rm aut}/K_{\rm eq}$, with
$K_{\rm eq}=(Y/(N_{\rm prod}A))^{1/\gamma}$, compares the individual
scale without credit to the observed productive scale. Here $Y$ is
the total production of a step and $N_{\rm prod}$ the number of
entities present in the productive phase, which can differ from the
surviving headcount at the end of the step. $K_{\rm eq}$ is the
fictitious stock that would give a representative entity the mean
production $Y/N_{\rm prod}$; for $\gamma\ne1$, this is not the mean
stock. A large $T$ therefore indicates a productively equivalent
stock far below the autarkic scale. This tension is useful in
particular at fixed $\delta$; it is neither a universal state
variable, nor another name for the whole
credit--mortality--population chain.

The coupling of distinct physical and economic spheres joins an
existing thermodynamic motivation \cite{herbert}, without reusing its
model or identifying money with exergy.

\subsection{Credit as cooperation, and the rate as sharing}
\label{sec:credit}

For two entities with the same concave technology, the marginal return
$f'(K)=A\gamma K^{\gamma-1}$ is larger for the less endowed one.
A transfer toward it therefore increases \emph{joint} production as
long as it does not exceed the equalization of marginal returns. This
result does not hold for any amount nor for any heterogeneous pair. It
is the sole motive for credit in this model: there is neither
anticipation, nor time preference, nor price.

Figure~\ref{fig:art_technologie} illustrates the role of technology and
the gain from lending. At fixed $\gamma$, $A$ multiplies extraction per
step; on its own it is not a power in watts. The exponent $\gamma$
measures individual returns to scale. Between zero and one, it makes
the function concave and allows a reallocation gain: calling it
directly a "propensity to collaborate" would be excessive. At
$\gamma=1$ and identical technology, this gain disappears; its
variation with $\gamma$ requires fixing the units and the states
compared.

\figurine{art_technologie}{1.0}{\textbf{Technology and cooperation gain.}
On the left, the technologies $f_1(K)=2K^{1/4}$ and $f_2(K)=K^{1/2}$
cross at $K=16$, for a production of 4 per step; the initial stocks
are distinct, $(K_1,K_2)=(20,4)$. The coefficients are expressed in a
single fixed stock unit. On the right, a loan of 1 toward 2 increases
joint production up to the equalization of marginal returns. The area
between gain and loss measures the joint surplus; the crossing of the
productions on the left is not this equalization.
These curves are analytical, at fixed states, without a statistical
interval.}

The market phase carries out $\lfloor\rho n\rfloor$ rounds for an
eligible pool of $n$ entities in the configurations studied; $\rho$ is
the market intensity (default value $1$). Encounters are pairwise.
The amount targets the equalization of marginal returns in the main
configuration. The free direction of the loan recovers the rich-to-poor
transfer in the homogeneous regime, not necessarily in the
heterogeneous regime.

There is therefore a \emph{cooperation surplus} to be shared, and the
interest rate is the sharing mechanism. For a contract of principal
$q$, let $L$ denote the production loss suffered by the giver by
parting with this stock, and $\Delta$ the joint production gain of the
pair -- the surplus. The service paid per step then decomposes as
\begin{equation}
r\,q \;=\; L \;+\; p\,\Delta.
\label{eq:partage}
\end{equation}
The number $p$ is the \emph{sharing}. The imposed \texttt{bargain} rule
explores $p\in[0,1]$; the implicit sharing of another rule can exceed
these bounds. At $p = 0$, the service exactly compensates the
productive loss computed at contract creation; the receiver keeps the
surplus. At $p = 1$, the lender also receives all of this surplus.
This instantaneous comparison does not guarantee the contract's
balance over its lifetime. Equation~\eqref{eq:partage} \emph{defines}
$p$ for any rate rule, which makes it possible to place a given rule on
the sharing scale -- this is what we exploit in
\S\ref{sec:partage}. The reference rule, used by default, sets $r$ to
the geometric mean of the marginal returns of the two entities.

\subsection{Bankruptcy and cascade}
\label{sec:faillite}

An entity goes bankrupt if it is in service default or if $NW_i < 0$.
The rule is deliberately brutal: its contracts are \emph{cancelled} --
its creditors lose the principal, with no recovery or buffer -- and its
residual capital is \emph{destroyed}. A default therefore directly
annihilates assets held by others and can render them insolvent in
turn; the new insolvencies form the next generation, and iteration
continues to the fixed point \emph{within the same step}.

This rule is what distinguishes a contagion from mere synchronization.
When a borrower dies, each lender loses the corresponding principal,
which defines a directed edge from the dead entity to the lender. An
\emph{avalanche} is a weakly connected component of this loss graph,
restricted to entities that died in the same step.
\begin{figure}[htbp]\centering
\begin{tikzpicture}[>=Stealth,font=\small,
 dead/.style={draw=red!65!black,fill=red!6,rounded corners,align=center,minimum width=23mm,minimum height=10mm},
 safe/.style={draw=green!45!black,fill=green!5,rounded corners,align=center,minimum width=26mm}]
\node[dead] (a) at (0,0) {$e_1$: root\\$NW=-1$};
\node[dead] (b) at (4,1) {$e_2$\\$2\to-2$};
\node[dead] (c) at (4,-1) {$e_3$\\$1\to-2$};
\node[dead] (d) at (8,1) {$e_4$\\$1\to-2$};
\node[safe] (e) at (8,-1.2) {$e_5$: survives\\$16\to10$};
\draw[->,thick] (a)--node[above,sloped] {loss $4$} (b);
\draw[->,thick] (a)--node[below,sloped] {loss $3$} (c);
\draw[->,thick] (b)--node[above] {loss $3$} (d);
\draw[->,dashed] (c)--node[below,sloped] {loss $3$} (e);
\draw[->,dashed] (d)--node[right,font=\scriptsize] {loss $3$} (e);
\end{tikzpicture}
\caption{\textbf{A directed propagation tree, not the direction of the loan.}
Example that is accounting-feasible with five entities. Before resolution,
their stocks are worth $(6,1,1,1,10)$; the loans are
$e_2\to e_1:4$, $e_3\to e_1:3$, $e_4\to e_2:3$,
$e_5\to e_3:3$ and $e_5\to e_4:3$.
The arrows in the diagram go, conversely, from the defaulting borrower to
the affected creditor. The component of the four deaths has one root and
two generations of propagation: $S=4$, $b_1=b_2=3/4$.
The losses suffered by $e_5$, a survivor, do not make it part of
the avalanche of deaths. The general model allows several parents,
unlike this tree-shaped example.}
\label{fig:cascade}
\end{figure}


Two propagation measures follow: the \emph{branching ratio}
$b_1 = 1 - (\text{roots}/\text{deaths})$, where a root is a bankruptcy
that entered directly in the step's initial queue, and the \emph{mean
offspring} $b_2$, the mean number of causal pairs between successive
generations per death. The two coincide
if and only if every non-root death has exactly one parent; their gap
measures the excess of edges per death:
$b_2-b_1=\sum_v(d_v-1)/N$,
with $d_v$ the number of previous-generation parents of a non-root and
$N$ the total number of deaths. This is neither the proportion of
victims with several parents, nor the number of individually
sufficient causes.
